% =====================================================================
% DocViz-Agent / QG-MDV — 실험 진행 계획서 v0.4.2 (한국어, Overleaf 호환)
% 본 문서는 학술 논문이 아닌 *내부 실험 진행 계획*임.
% 학술 논문 본문은 docviz_paper_draft_v0.4.2_ko.tex 참고.
% 작성 기준일: 2026-06-05
% =====================================================================

\documentclass[11pt]{article}

\usepackage[hangul]{kotex}
\usepackage[english]{babel}
\usepackage{amsmath,amssymb}
\usepackage{booktabs}
\usepackage{multirow}
\usepackage{xcolor}
\usepackage{hyperref}
\usepackage{geometry}
\geometry{margin=1in}
\usepackage{enumitem}
\usepackage[normalem]{ulem}
\usepackage{tcolorbox}

\title{DocViz-Agent v0.4.2 \\
       단계적 실험 진행 계획서 \\
       \large{(Internal Execution Plan, not for publication)}}
\author{연구진 (placeholder)}
\date{2026년 6월}

\begin{document}
\maketitle

\begin{abstract}
본 문서는 DocViz-Agent v0.4.2 연구의 실험 진행 계획서이다. 학술 논문 본문 (docviz\_paper\_draft\_v0.4.2\_ko.tex)과 분리되어 \emph{내부 실험 운영 용도}로 작성되며, 다섯 개의 단계와 각 단계의 검증 기준 (진행 결정 관문, go/no-go gate), 누적 비용 추정, 위험 차단 전략을 정의한다. 단계 분리의 목적은 v0.4.1 시제품 (prototype) 단계에서 발견된 평가 결함 (eval artifact) 같은 위험을, 전체 비용을 투입하기 전에 단계별로 차단하는 데 있다.
\end{abstract}

% =====================================================================
\section{문서 목적 및 위치}
% =====================================================================

본 문서는 학술 논문 본문이 아니다. 학술 논문에는 일반적으로 단계별 작업 계획이나 진행 결정 관문 같은 운영 항목이 포함되지 않으므로, 본 연구는 두 문서를 분리하여 운영한다.

\begin{itemize}[leftmargin=*]
    \item \textbf{학술 논문 본문} (\texttt{docviz\_paper\_draft\_v0.4.2\_ko.tex}): 외부 학회 제출용. 서론, 관련 연구, 과제 정형화, 방법, 평가 프레임워크, 실험 설정, 결과, 절제 분석, 검증, 논의의 표준 학술 논문 구조.
    \item \textbf{실험 진행 계획서} (\texttt{docviz\_execution\_plan\_v0.4.2\_ko.tex}, 본 문서): 연구진 내부용. 5단계 진행 계획, 단계별 검증 기준, 누적 비용, 위험 차단 전략.
\end{itemize}

두 문서의 결정 사항은 일관되며, 학술 논문의 § 6 ``실험 설정'' 절에 본 진행 계획의 일부 결과가 *최종 측정 결과*로만 보고된다.

% =====================================================================
\section{단계 분리의 근거}
% =====================================================================

v0.4.1 시제품 단계에서 다음 문제가 발견되었다. 자체 정확 일치 (exact match) 셀 F1과 자체 정규식 (regex) 기반 Mermaid 파서를 도입한 결과, 모든 비교 대상의 점수가 97\% 영점에 수렴하는 바닥 상태 (floor state)가 관찰되었다. 이 사실이 확인된 시점에 이미 다수의 다중 시드 (multi-seed) 측정과 여러 백본에 대한 측정이 진행된 후였으며, 평가 파이프라인의 결함이 *전체 비용 투입 후*에야 발견되었다는 점이 결정적 손실이었다.

본 v0.4.2는 동일한 종류의 위험을 단계별로 사전 차단한다. 가장 위험한 단계 — 평가 파이프라인 자체의 변별력 (discriminative power) 검증 — 을 최소 비용으로 가장 먼저 수행하며, 통과한 경우에만 다음 단계로 진입한다. 단계별 누적 비용은 점진적으로 증가하므로, 어느 단계에서 위험이 발견되더라도 그 단계까지의 누적 비용만 손실되고 다음 단계의 큰 비용은 차단된다.

% =====================================================================
\section{단계별 진행 계획}
% =====================================================================

\subsection{1단계 — 기반 사전 검증 (Foundation Pilot)}

\textbf{기간 / 비용}: 1주 / 약 5달러

\textbf{목표}: 평가 파이프라인이 변별력 있는 점수를 산출하는지 확인. 시제품 단계의 바닥 상태 재발을 차단.

\textbf{작업}: Loong 벤치마크에서 30개 표본을 추출하여 Qwen3.5-397B 단일 백본으로 두 파이프라인 (B6 완전판, B5 단일 시도 LLM)을 실행한다. 모든 다섯 종 mermaid 하위 유형에서 렌더링-시각 언어 모델 (VLM) 추출 파이프라인을 검증하며, 차트 산출물에 대해 Doc2Chart의 CHARTEVAL 귀속 평가를 적용한다.

\textbf{검증 기준 (Go gate)}:
\begin{itemize}[leftmargin=*]
    \item Mermaid 렌더링 → VLM 추출 성공률이 다섯 종 하위 유형 \emph{모두}에서 85\% 이상
    \item 차트 귀속 성공률 85\% 이상
    \item B6와 B5의 점수 차이가 적어도 두 지표에서 $+0.020$ 이상 (변별력 확인)
    \item 3시드 분산이 5\% 이하 (재현성 확인)
\end{itemize}

\textbf{미통과 시}: 평가 파이프라인 수정 후 1단계 반복. \emph{다음 단계 진입 금지}. 시제품 단계의 바닥 상태 재발이 의심되는 경우 이 단계에서 발견되어야 한다.

\subsection{2단계 — 최소 실효 실험 (Minimal Viable Experiment)}

\textbf{기간 / 비용}: 1.5주 / 약 20달러

\textbf{목표}: 제안 방법 (B6 DocViz-Agent)의 실제 효과 입증. 1단계가 평가 파이프라인의 신뢰성을 확인한 후, 2단계는 본 연구의 방법 자체가 효과를 보이는지를 검증한다.

\textbf{작업}: Loong에서 100개 표본 (다섯 추론 난이도 유형 각 20개)을 추출하여 Qwen3.5-397B 단일 백본으로 세 비교 대상 (B5, B7 SelfRefine, B6)을 실행한다. 모든 핵심 지표 — DiagramEval 노드/경로 F1, Doc2Chart CHARTEVAL, Evidence F1, Intent Coverage, M1 렌더링 성공률, M5 CLIPScore — 를 측정한다. 동시에 Prolific에서 30개 표본에 대한 자연스러움 점검을 수행한다 (약 30달러).

\textbf{검증 기준}:
\begin{itemize}[leftmargin=*]
    \item 강 신호 (strict gate) 기준: B6가 두 개 이상 지표에서 가장 강한 비교 대상 대비 $+0.020$ 이상
    \item B6와 B7의 차이가 3시드 표준편차 기준으로 통계적으로 유의
\end{itemize}

\textbf{미통과 시}: B6 파이프라인 (CIS / TMG / SAO 세 축) 보완 후 2단계 반복. 시제품 단계 측정에서 TMG 단독 절제 시 가장 큰 격차가 관찰되었으므로, 미통과 분석은 TMG 축의 운영 방식부터 점검한다.

\subsection{3단계 — 다중 출처 강건성 (Multi-source Robustness)}

\textbf{기간 / 비용}: 2주 / 약 30달러

\textbf{목표}: Loong 외 출처 벤치마크로의 일반화 (generalization). 2단계가 단일 출처에서의 효과를 입증한 후, 3단계는 그 효과가 출처에 무관하게 재현되는지 검증한다.

\textbf{작업}: 2단계 설정에 DocHop-QA와 FinMMDocR을 추가하여 세 출처에서 각 60개 표본씩을 측정한다. 백본은 여전히 Qwen3.5-397B로 단일 유지하여 출처 효과를 분리한다.

\textbf{검증 기준}:
\begin{itemize}[leftmargin=*]
    \item B6의 우위가 세 출처 중 두 개 이상에서 재현
    \item 출처 간 방법 효과의 분산이 50\% 이하 (단일 출처가 결과를 주도하지 않음)
\end{itemize}

\textbf{미통과 시}: 특정 출처에서만 작동한다는 신호로 해석되며, 질의 생성 파이프라인 (\S\ref{sec:querygen})의 출처별 프롬프트 적응을 점검한다.

\subsection{4단계 — 백본 횡단 검증과 진단적 발견 (Cross-backbone, C4 Finding)}

\textbf{기간 / 비용}: 2주 / 약 60달러 (클로즈드 API)

\textbf{목표}: 본 연구의 네 번째 기여인 다문서 출처 귀속 격차 (multi-document grounding gap)를 네 개의 백본에서 입증.

\textbf{작업}: 3단계 설정 (3 출처 × 60 표본 × 3 비교 대상)을 네 개 백본 — Qwen3.5-397B, DeepSeek-V4-Flash, GPT-5-mini, Claude Opus 4.8 — 으로 확장하여 측정한다.

\textbf{검증 기준}:
\begin{itemize}[leftmargin=*]
    \item 진단적 발견 (C4): Evidence F1의 격차가 세 개 이상의 백본에서 일관되게 관찰됨
    \item B6의 우위가 세 개 이상의 백본에서 재현
\end{itemize}

\textbf{미통과 시}: 백본별 실패 양상을 분석한다. 특정 백본에서 시각화 생성 자체가 작동하지 않는 경우는 그 백본을 비교 풀에서 교체할 수 있다.

\subsection{5단계 — 완전 말뭉치와 외부 보류 평가 (Full Corpus, Held-out)}

\textbf{기간 / 비용}: 2주 / 약 100달러

\textbf{목표}: 300개 표본 전체와 세 개의 외부 보류 평가 벤치마크에 대한 측정으로 논문 본 결과 (main result) 완성.

\textbf{작업}: 일곱 개의 출처 벤치마크 모두 (총 300 표본)에 네 백본과 일곱 비교 대상을 적용한다. Text2Vis, Doc2Chart, SciDoc2DiagramBench의 세 외부 보류 평가 벤치마크를 추가로 측정한다. 모든 측정은 시드 42, 43, 44의 3시드 평균과 표준편차로 보고한다.

\textbf{검증 기준}:
\begin{itemize}[leftmargin=*]
    \item 강 신호 기준이 완전 말뭉치에서 통과
    \item B6가 다섯 추론 난이도 유형 중 네 개 이상에서 우위
    \item 외부 보류 평가에서 전문가 모델 (specialist) 대비 $-7\%$ 포인트 이내
\end{itemize}

\textbf{미통과 시}: 결과 시나리오 분석 (\S\ref{sec:scenarios})에 따라 논문 위치 (main vs Findings)를 조정한다.

\subsection{6단계 — 이미지 평가, 사람 검증, 논문 작성}

\textbf{기간 / 비용}: 2주 / 약 50달러 + Prolific 비용

\textbf{목표}: 이미지 단계 평가 완료, 사람 정합성 검증, 논문 그림과 표 완성.

\textbf{작업}: M5 CLIPScore를 모든 시각화 산출물에 적용한다 (결정론적, 비용 0). A5 시각 리커트 척도를 Claude Sonnet 4.6 비전 API 배치 모드로 100개 부집합에 적용한다 (약 11달러). Prolific에서 50개 시각화에 대한 세 명의 평가자가 자연스러움을 평가하는 정박 (anchor) 작업을 수행한다. 역방향 질의응답 (reverse question answering) L4 정박을 100개 질의 × 7 비교 대상에 적용한다 (약 3달러).

\textbf{검증 기준}:
\begin{itemize}[leftmargin=*]
    \item 사람 점검의 스피어만 상관 $r \geq 0.4$ (DiagramEval 원논문 보고값 0.43의 범위)
    \item 모든 논문 표의 빈 칸이 채워짐
\end{itemize}

% =====================================================================
\section{누적 비용 및 시간}
% =====================================================================

\begin{table}[h]
\centering
\small
\begin{tabular}{lrrl}
\toprule
단계 & 누적 비용 (달러) & 누적 시간 & 차단되는 위험 \\
\midrule
1단계 사전 검증 & 5 & 1주 & 평가 결함, 바닥 상태 \\
2단계 최소 실효 & 25 & 2.5주 & 방법의 실재성 \\
3단계 다중 출처 & 55 & 4.5주 & 단일 출처 의존 \\
4단계 백본 횡단 & 115 & 6.5주 & 진단적 발견의 검증 \\
5단계 완전 말뭉치 & 215 & 8.5주 & 본 결과의 일반성 \\
6단계 검증 작성 & 315 (+ Prolific) & 10.5주 & 심사위원 정박 (reviewer anchor) \\
\bottomrule
\end{tabular}
\caption{단계별 누적 비용과 누적 시간. 각 단계에서 미통과 판정 시 그 단계까지의 누적 비용만 손실되며, 다음 단계의 큰 비용은 차단된다.}
\end{table}

% =====================================================================
\section{질의 생성 파이프라인 상세}
\label{sec:querygen}
% =====================================================================

질의 생성은 세 단계 파이프라인으로 수행되며, Loong, Doc2Chart, Text2Chart31의 세 선행 연구에서 검증된 메커니즘을 차용한다. 모든 출처 벤치마크에 동일한 프롬프트 템플릿이 적용되며, 출처별 추론 표지만 슬롯에 삽입된다.

첫째 단계는 압축 문서 표현 (compressed doc representation)이다. Loong의 자기 보고된 자료 생성 절차 (논문 §3.3.2)에 따라 각 문서를 핵심 사실 글머리 목록으로 압축한다. 원문 텍스트를 그대로 사용하는 대신 압축된 표현을 사용함으로써 후속 단계의 모델이 핵심 정보에 집중할 수 있도록 한다. 압축에는 GPT-4o-mini 모델이 사용된다.

둘째 단계는 과제별 프롬프트와 출력 유형 휴리스틱 (task-specific prompt with output-type heuristic)이다. Loong의 네 가지 과제 설명 방식과 Doc2Chart의 출력 유형 휴리스틱을 결합하여, 추론 난이도 유형 × 출력 구조 조합별로 25개의 프롬프트 템플릿을 준비한다. 각 템플릿은 한 차례의 학습 시연 예시를 포함하며, 출처 벤치마크의 표본 단위 추론 표지가 정박 근거로 명시된다.

셋째 단계는 자기 평가 관문 (self-evaluation gate)이다. Text2Chart31의 메커니즘 (논문 §3.3)을 차용하여, 동일 모델이 생성된 질의를 다문서 필요성, 시각화 적합성, 사용자 의도 현실성의 세 기준에서 평가하며, 세 기준을 모두 충족하는 경우에만 채택한다. 한 차례의 재생성 (retry)이 허용된다.

전체 파이프라인의 비용은 300개 질의에 대해 LLM 호출 900회 (단계당 1회씩) × GPT-4o-mini 단가 약 0.001달러로 총 약 0.96달러이다.

% =====================================================================
\section{난이도 필터링 상세}
% =====================================================================

질의 생성 후 난이도 필터링이 수행된다. 평가 풀과 분리된 두 개의 프론티어 모델 — GPT-5 완전판과 Claude Opus 4.7 — 이 다수결로 사용된다. 두 모델 모두 본 연구의 평가 풀 (Qwen3.5-397B, DeepSeek-V4-Flash, GPT-5-mini, Claude Opus 4.8)에 포함되지 않으므로 순환 평가 (circular evaluation)의 위험이 차단된다.

두 모델이 모두 통과하는 질의는 너무 쉬운 것으로 간주되어 제거되며, 두 모델 모두 실패하는 질의 100개를 엄격 표본 (strict)으로, 한 모델만 실패하는 질의 200개를 완화 표본 (soft)으로 구성하여 총 300개의 다문서 묶음을 확보한다. 비용은 600개의 후보 질의에 대해 두 모델 × 약 0.06달러로 약 132달러 일회성이다.

% =====================================================================
\section{골드 부집합 사람 검증 상세}
% =====================================================================

300개 표본 중 90개의 부집합 (challenge type 각 15개, 도메인 상한 6 적용)을 Prolific에서 사람 검증한다. 부집합 선정은 시드 42의 계층적 무작위 표집 (stratified random sampling)으로 수행된다.

검증은 두 단계로 구성된다. 다중 LLM 합집합 (multi-LLM union) 추출 단계에서는 GPT-5-mini와 Claude Opus 4.8이 동일 프롬프트로 독립적으로 증거 범위 (evidence span), 사실 (fact), 표, 그래프, 의도를 추출하며, 두 결과의 합집합을 취한 후 중복을 제거한다. Prolific 사람 검증 단계에서는 90개 표본 각각에 대해 세 명의 평가자가 두 가지 과제를 수행한다.

출처 정확성 (source correctness) 과제에서는 각 추출된 사실이 인용된 원문 범위에 실제로 존재하는지를 이진 (binary) 판정하며, 평가자 세 명 중 2/3 이상이 동의할 때 해당 사실을 보존한다. 완전성 (completeness) 과제에서는 평가자가 자유 텍스트로 누락된 중요 사실을 추가하며, 추가된 사실은 네 번째 평가자가 재검증한 후 합병된다.

시제품 단계에서 운영하였던 과추출 (over-extraction) 점검 과제는 셀 수준 (cell-level) 골드를 본 연구에서 자체 구현하지 않으므로 폐기된다. 총 사람 검증 비용은 약 200달러로 추정된다.

% =====================================================================
\section{결과 시나리오와 예비 논문 계획}
\label{sec:scenarios}
% =====================================================================

5단계 완료 시점의 결과에 따라 논문의 위치와 강조점이 결정된다. 본 절은 시나리오별 예비 계획을 정의한다.

\textbf{최선 시나리오}: 외부 보류 평가에서 전문가 대비 $-5\%$ 포인트 이내, QG-MDV에서 강 신호 $+10\%$ 포인트, 교차 과제 평균 $+12\%$ 포인트, C4 발견 강함. → 주류 학회 (EMNLP main).

\textbf{양호 시나리오}: $-7\%$ 이내, $+8\%$, $+8\%$, C4 강함. → 주류 또는 Findings.

\textbf{보통 시나리오}: $-10\%$ 이내, $+5\%$, $+5\%$, C4 약함. → Findings.

\textbf{B6 효과 약함}: 외부 보류 평가는 양호하나 QG-MDV에서 $+3\sim5\%$. → 자원 (resource) 논문으로 Findings 제출 가능 (벤치마크 자체가 기여).

\textbf{외부 보류 평가 실패}: $-15\%$ 이상의 격차. → 분석 (analysis) 논문으로 재구성. 전문화 대 일반화 (specialization vs generalization) 트레이드오프 분석.

\textbf{진단적 발견만 유의}: B6 효과는 약하나 C4 발견은 강함. → 발견 (finding) 논문으로 재구성. 다문서 출처 귀속 격차가 본 논문의 주된 기여가 된다.

모든 시나리오에서 출판 가능 (publishable)하다는 점이 본 연구 설계의 강점이다.

% =====================================================================
\section{위험 차단 전략}
% =====================================================================

각 단계의 검증 기준이 미통과되는 경우의 대응을 명시한다.

\textbf{1단계 미통과}: 평가 파이프라인 자체의 결함이 의심된다. 시제품 단계의 자체 셀 F1과 자체 파서를 폐기한 결정의 정당성이 흔들리는 신호이므로, Doc2Chart의 CHARTEVAL과 DiagramEval의 원본 코드를 정확히 따르고 있는지 점검한다. 자체 추가 처리가 들어있지 않은지 코드 검토.

\textbf{2단계 미통과}: B6 파이프라인이 단일 시도 LLM (B5) 대비 변별 가능한 우위를 보이지 않는다. CIS, TMG, SAO 세 축의 운영 방식을 점검한다. 시제품 단계에서 TMG 단독 절제 시 가장 큰 격차가 관찰되었으므로, TMG의 십개 시각화 기본형 풀 노출이 정상 작동하는지부터 점검한다.

\textbf{3단계 미통과}: 단일 출처 (Loong) 효과가 다른 출처로 확장되지 않는다. 질의 생성 파이프라인의 출처별 적응 (compressed doc representation의 출처별 글머리 형식, task-specific prompt의 출처별 표지 매핑)을 점검한다.

\textbf{4단계 미통과}: 백본 종속성이 강하다. 특정 백본에서 시각화 생성이 본질적으로 작동하지 않는 경우, 그 백본을 평가 풀에서 교체한다. 단 클로즈드 프론티어 (Opus 4.8)는 비교 풀에서 제외할 수 없으므로, 미통과 시 진단적 발견 (C4)의 주장 범위를 ``오픈 백본에 한정''으로 축소한다.

\textbf{5단계 미통과}: 완전 말뭉치에서 강 신호 기준이 충족되지 않는다. 결과 시나리오 분석 (\S\ref{sec:scenarios})에 따라 논문 위치를 Findings로 조정하거나, 발견 논문으로 재구성한다.

% =====================================================================
\section{시제품 단계에서 폐기된 항목 정리}
% =====================================================================

v0.4.2 진행 과정에서 시제품 단계 또는 초기 설계에서 검토되었으나 폐기된 항목을 정리한다. 폐기 결정의 근거는 효율성, 위험성, 정당화 부담의 세 측면에서 평가되었다.

\begin{table}[h]
\centering
\small
\begin{tabular}{lp{7.5cm}}
\toprule
폐기 항목 & 폐기 근거 \\
\midrule
시제품 단계의 네 축 RocketEval 평가 (scope-v3) & Doc2Chart $r=0.71$과 DiagramEval $r=0.43$의 발표된 사람 정합성으로 충분히 검증 인계 가능. 자체 평가 설계의 정당화 부담 (오염 방지 규칙 등) 회피. \\
ViviBench 외부 보류 평가 & 공식 평가 코드 미공개로 자체 재구현 시 재현성 위험. ``재구현이 발표 지표와 동일한가'' 공격에 방어 없음. \\
Plot2Code 필수 보류 평가 지정 & 본 연구 설정과 정합성 약함. 선택적 부록 민감도 분석으로 강등. \\
FSOS 가중치 합산 점수 & 자체 가중치 결정의 정당화 부담. 네 개 지표의 병렬 보고로 대체. \\
자체 정상화 계층 (정상화: 엣지 관계 군집화, 수치 단위 정규화, 범위 중첩 완화) & 자체 규칙 기반 구현의 결함 위험. 선행 연구의 공식 평가 코드 그대로 사용하여 검증 인계. \\
\texttt{claude -p} 명령줄 도구 (CLI) 사용 경로 & 사용량 제한 (rate limit)으로 2,800회 배치 호출 차단됨. Anthropic API SDK + 배치 모드로 전환. \\
HotpotQA / 2WikiMultihopQA / MuSiQue 등 Wikipedia 기반 데이터 & 본 연구의 ``실세계 원문 자료'' 정의에 부적합. \\
\bottomrule
\end{tabular}
\caption{v0.4.2에서 폐기된 항목과 폐기 근거. 각 폐기 결정은 효율성, 위험성, 정당화 부담의 세 측면에서 평가되었다.}
\end{table}

% =====================================================================
\section{문서 간 일관성 점검}
% =====================================================================

본 진행 계획서와 학술 논문 본문 (\texttt{docviz\_paper\_draft\_v0.4.2\_ko.tex}) 사이의 일관성은 다음 항목에서 유지된다.

학술 논문 본문의 ``실험 설정'' 절에 본 진행 계획의 다음 항목이 반영된다: 백본 풀 구성 (네 모델 등급 분산), 말뭉치 구성 (일곱 출처 벤치마크 정박), 질의 생성 파이프라인 (세 단계 구조), 골드 부집합 사람 검증 (90개 표본).

학술 논문 본문에 반영되지 \emph{않는} 본 진행 계획의 항목은 다음과 같다: 단계별 작업 일정과 비용, 단계별 검증 기준 (진행 결정 관문), 위험 차단 전략, 결과 시나리오와 예비 논문 계획, 폐기 항목 정리. 이러한 항목들은 학술 논문 본문의 자리가 아니므로 본 계획서에만 포함된다.

\end{document}
